\subsubsection{Quantum Ising Model}\label{sec:quantum-ising-model}

In the previous section, we studied the classical \textbf{Ising model}, where each spin is represented by a classical variable:
\begin{equation*}
    s_i \in \{-1, +1\} \quad \forall i \in \{1, \dots, n\}
\end{equation*}
Where $n$ is the number of spins in the system. The energy of the system was:
\begin{equation*}
    H_{\text{Ising}} = \sum_i h_i s_i + \sum_{i > j} J_{ij} s_i s_j
\end{equation*}
This equation assigns an energy to each configuration of the spins. However, the spins are still classical objects that can only be in one of two possible states.

\highspace
\begin{flushleft}
    \textcolor{Green3}{\faIcon{atom} \textbf{From Classical Spins to Quantum Spins}}
\end{flushleft}
To obtain a quantum model, we replace the classical spin variables $s_i$ with quantum operators. But the natural question is: \textbf{\emph{which quantum operators should we use to represent the spins?}}

\highspace
We want something that behaves like a spin. Since a spin can be in two states:
\begin{equation*}
    s_i \in \{-1, +1\}
\end{equation*}
We look for a quantum operator whose eigenvalues are also $\{-1, +1\}$. The Pauli $Z$ operator is a perfect candidate for this:
\begin{equation*}
    Z = \begin{bmatrix}
        1 & 0 \\
        0 & -1
    \end{bmatrix}
\end{equation*}
Because its eigenvectors are $\ket{0}$ and $\ket{1}$:
\begin{equation*}
    \ket{0} = \begin{bmatrix}
        1 \\
        0
    \end{bmatrix}, \quad \ket{1} = \begin{bmatrix}
        0 \\
        1
    \end{bmatrix}
\end{equation*}
And its eigenvalues are $+1$ and $-1$:
\begin{equation*}
    Z \ket{0} = +1 \ket{0}, \quad Z \ket{1} = -1 \ket{1}
\end{equation*}
This means that we can represent the spin $s_i$ as the operator $Z_i$, where $Z_i$ acts on the $i$-th qubit:
\begin{center}
    \begin{tabular}{@{} c c @{}}
        \toprule
        \textbf{Classical Spin} & \textbf{Quantum Spin} \\
        \midrule
        $+1$ & $\ket{0}$ \\
        $-1$ & $\ket{1}$ \\
        \bottomrule
    \end{tabular}
\end{center}

\noindent
This mapping allows us to rewrite the classical Ising Hamiltonian in terms of quantum operators:
\begin{equation}\label{eq:quantum-ising-hamiltonian}
    \hat{H}_{\text{Ising}} = \sum_i h_i Z_i + \sum_{i > j} J_{ij} Z_i Z_j
\end{equation}
The structure is identical to the classical case, the only difference is that the \hl{classical variables have been replaced by quantum operators.}